\PassOptionsToPackage{table}{xcolor}
\documentclass[aspectratio=169,11pt,t]{beamer}
\usepackage[russian]{academicbeamer}
\setshorttitle{BelOAI~\textbar~Лекция 1}

\usetikzlibrary{arrows.meta}
\graphicspath{{assets/}}

% Deck-local styling for the generated charts (assets/*.tex) and flow diagrams.
\definecolor{acOrange}{RGB}{225,115,45}
\colorlet{acAccent}{acMid}
\colorlet{acInk}{acDark}
\colorlet{acCard}{white}
\colorlet{acText}{black!85}
\tikzset{
  chart axis/.style={draw=acGray!70, line width=0.6pt},
  chart grid/.style={draw=acRule, line width=0.4pt},
  chart label/.style={font=\scriptsize, text=acGray, inner sep=2pt},
  chart curve/.style={draw=acMid, line width=1.2pt, line cap=round},
  chart train/.style={circle, fill=acOrange, inner sep=0pt, minimum size=4.6pt},
  chart test/.style={circle, draw=black!80, fill=white, line width=0.75pt,
    inner sep=0pt, minimum size=4.6pt}
}

% Zebra tables in the theme palette: \zebra{first-shaded-row} before a tabular.
\newcommand{\zebra}[1]{\rowcolors{#1}{acLight}{white}\renewcommand{\arraystretch}{1.45}\setlength{\tabcolsep}{7pt}\normalsize}

% Numbered rows: \begin{numrows} 1 & text\\ ... \end{numrows}
\newenvironment{numrows}{\zebra{1}\begin{tabular}{>{\bfseries\large\color{acMid}}c >{\raggedright\arraybackslash}p{\dimexpr\linewidth-1.7cm\relax}}}{\end{tabular}}

% Ragged-right text: no stretched word spacing or ugly hyphenation in narrow blocks.
\usepackage{ragged2e}
\addtobeamertemplate{block begin}{}{\advance\rightskip 0pt plus 4em\relax}
\setbeamerfont{block title}{size=\normalsize,series=\bfseries}
\setbeamerfont{block body}{size=\normalsize}
\hyphenpenalty=10000
\exhyphenpenalty=10000

\title{Машинное обучение и где оно обитает}
\subtitle{Задачи, метрики качества, данные, модели и обучение}
\author{Автор презентации: Павел Дылевский\\
  Доработал и обновил: Алексей Чернуха}
\date{По статье 1.3 «Машинное обучение» из учебника по ML (Яндекс)}

\begin{document}

\begin{frame}[plain]
  \vspace{1.1cm}
  \hspace{0.02\paperwidth}\begin{minipage}{0.9\paperwidth}
  {\color{acRule}\hrule height 0.9pt}\vspace{0.8em}
  {\usebeamerfont{title}\usebeamercolor[fg]{title}\inserttitle\par}
  \vspace{0.7em}
  {\large\color{acMid}\insertsubtitle\par}
  \vspace{0.8em}{\color{acRule}\hrule height 0.9pt}
  \vspace{1.4em}
  {\small\color{acGray}\insertauthor\par}
  \vspace{0.9em}
  {\small\color{acGray}\insertdate}
  \end{minipage}
\end{frame}

\section{Введение}

\begin{frame}[c]{Зачем нужно машинное обучение?}
  \lead{Вопрос на подумать.}
  Алан Тьюринг расшифровывал Энигму перебором вариантов --- машина справилась быстрее людей.

  \bigskip
  \begin{block}{Есть задачи, которые человеку даются легко, а программой их почти не записать}
    Назовите примеры таких задач.
  \end{block}
\end{frame}

\begin{frame}[c]{Ответ: легко человеку, трудно запрограммировать}
  \begin{columns}[T]
    \begin{column}{0.52\textwidth}
      \begin{numrows}
        1 & Перевести текст с языка на язык\\
        2 & Поставить диагноз по симптомам\\
        3 & Выбрать, какой из двух документов лучше подходит запросу\\
        4 & Сказать, что изображено на картинке\\
        5 & Оценить, по какой цене продастся квартира\\
      \end{numrows}\par
    \end{column}
    \begin{column}{0.44\textwidth}
      \vspace*{-4pt}
      \begin{block}{Что у них общего}
        \begin{enumerate}\setlength\itemsep{0.3em}
          \item Решение --- функция из объекта в предсказание.
          \item Единственно верного ответа нет: переводы бывают разными, врачи ошибаются.
          \item Правильных ответов много, а неправильные легко придумать.
        \end{enumerate}
      \end{block}
    \end{column}
  \end{columns}
\end{frame}

\begin{frame}[c]{Что такое машинное обучение?}
  \vfill
  \lead{Вопрос.}

  \medskip
  {\Large\bfseries\color{acDark} Объясните своими словами, как человеку, далёкому от IT.\par}
  \vfill
\end{frame}

\begin{frame}[c]{Определение машинного обучения}
  \lead{Компьютер сам находит закономерности в данных.}
  Вместо того чтобы человек вручную прописывал правила, алгоритм выучивает зависимость между
  объектами и ответами по примерам и применяет её к новым объектам.

  \medskip
  \begin{columns}[T]
    \begin{column}{0.31\textwidth}
      \begin{block}{Данные}
        Примеры: объекты и правильные ответы к ним.
      \end{block}
    \end{column}
    \begin{column}{0.31\textwidth}
      \begin{block}{Модель}
        Функция $f(x) \to y$: по объекту $x$ выдаёт предсказание $y$.
      \end{block}
    \end{column}
    \begin{column}{0.31\textwidth}
      \begin{block}{Обучение}
        Подбор параметров модели так, чтобы ответы на данных были как можно лучше.
      \end{block}
    \end{column}
  \end{columns}
\end{frame}

\begin{frame}[c]{Из чего состоит задача}
  \centering
  \begin{tikzpicture}[
      box/.style={rectangle,minimum width=0.27\textwidth,minimum height=0.9cm,
        text=white,font=\bfseries,align=center},
      arr/.style={-{Stealth[length=2.5mm]},line width=1.2pt,color=acGray}]
    \node[box,fill=acDark] (x) {Объект $x$};
    \node[box,fill=acMid,right=0.8cm of x] (f) {Модель $f$};
    \node[box,fill=acDark,right=0.8cm of f] (y) {Предсказание $y$};
    \draw[arr] (x) -- (f);
    \draw[arr] (f) -- (y);
  \end{tikzpicture}

  \medskip
  \begin{columns}[T]
    \begin{column}{0.31\textwidth}
      \begin{block}{Объекты и таргеты}
        Объект --- картинка, история болезни, пользователь сервиса.
        Таргет --- ответ: подпись, диагноз, уход пользователя.
      \end{block}
    \end{column}
    \begin{column}{0.31\textwidth}
      \begin{block}{Модель}
        Функция $f(x) \to y$. Принимает признаки объекта $x$ и выдаёт предсказание $y$.
      \end{block}
    \end{column}
    \begin{column}{0.31\textwidth}
      \begin{block}{Датасет}
        Обучающая выборка: объекты вместе с известными ответами. По ней учится модель.
      \end{block}
    \end{column}
  \end{columns}
\end{frame}

\section{Постановка задачи}

\begin{frame}[c]{Три парадигмы обучения}
  \begin{columns}[T]
    \begin{column}{0.31\textwidth}
      \begin{block}{С учителем}
        Для каждого объекта известен правильный ответ. Учимся предсказывать ответ для новых объектов.

        \smallskip
        \lead{Пример:} диагноз по симптомам.
      \end{block}
    \end{column}
    \begin{column}{0.31\textwidth}
      \begin{block}{Без учителя}
        Есть только данные, ответов нет. Ищем структуру: например, группы похожих объектов.

        \smallskip
        \lead{Пример:} группировка новостей по сюжетам.
      \end{block}
    \end{column}
    \begin{column}{0.31\textwidth}
      \begin{block}{С подкреплением}
        Агент действует в среде и получает награду за удачные действия.

        \smallskip
        \lead{Пример:} робот учится запрыгивать на стол.
      \end{block}
    \end{column}
  \end{columns}
\end{frame}

\begin{frame}[c]{Что предсказывает модель?}
  \lead{Вопрос на подумать.}
  \begin{numrows}
        1 & Время поездки на каршеринге\\
        2 & Кликнет ли пользователь по рекламе\\
        3 & Тема научной статьи: математика, биология, психология\ldots\\
        4 & Теги ресторана: «кофе», «веранда», «живая музыка»\\
        5 & Порядок документов в поисковой выдаче\\
      \end{numrows}\par
  \medskip
  \textbf{Чем различаются ответы? Какие группы задач вы видите?}
\end{frame}

\begin{frame}[c]{Обучение с учителем: виды задач}
  \small
  \zebra{2}\small
  \begin{tabular}{>{\raggedright\arraybackslash}p{0.286\textwidth}>{\raggedright\arraybackslash}p{0.593\textwidth}}
    \rowcolor{acDark} \textcolor{white}{\textbf{Задача}} & \textcolor{white}{\textbf{Что предсказываем}}\\
    \textcolor{acMid}{\textbf{Регрессия}} & Таргет --- число или вектор. Время поездки на каршеринге, спрос на товар.\\
    \textcolor{acMid}{\textbf{Бинарная классификация}} & Да или нет. Кликнет ли по рекламе, вернёт ли кредит.\\
    \textcolor{acMid}{\textbf{Многоклассовая}} & Один класс из нескольких. Тема научной статьи.\\
    \textcolor{acMid}{\textbf{Многометочная}} & Несколько классов сразу, они могут пересекаться. Теги ресторана.\\
    \textcolor{acMid}{\textbf{Ранжирование}} & Упорядочить объекты. Поисковая выдача по релевантности.\\
    \textcolor{acMid}{\textbf{Сложные ответы}} & Сегментация изображения, перевод текста, генерация новых объектов.\\
  \end{tabular}
\end{frame}

\begin{frame}[c]{Определите тип задачи}
  \lead{Вопрос на подумать.}
  \begin{numrows}
        1 & Предсказать курс евро к доллару на следующий день\\
        2 & Стилизовать текст: «Пиппина и Мерри похитили!» $\to$ «Граждане Тук и Брендибак были похищены неустановленными лицами»\\
        3 & Определить, котик или собачка на картинке\\
        4 & Научить робокота запрыгивать на стол из любой позы\\
        5 & Найти наборы товаров, которые покупают вместе\\
      \end{numrows}\par
\end{frame}

\begin{frame}[c]{Ответы: тип задачи}
  \small
  \zebra{2}
  \begin{tabular}{>{\raggedright\arraybackslash}p{0.267\textwidth}>{\raggedright\arraybackslash}p{0.611\textwidth}}
    \rowcolor{acDark} \textcolor{white}{\textbf{Задача}} & \textcolor{white}{\textbf{Тип}}\\
    1. Курс евро & Регрессия: предсказываем число\\
    2. Стилизация текста & Генерация новых объектов на основе существующих\\
    3. Котик или собачка & Бинарная классификация: два класса\\
    4. Робокот & Обучение с подкреплением: перебирать траектории и узнавать ответ слишком долго и дорого\\
    5. Наборы товаров & Обучение без учителя: ответов нет\\
  \end{tabular}
\end{frame}

\section{Критерии качества}

\begin{frame}[c]{Как понять, что модель хороша}
  \small
  \zebra{1}
  \begin{tabular}{>{\raggedright\arraybackslash}p{0.191\textwidth}>{\raggedright\arraybackslash}p{0.688\textwidth}}
    \textcolor{acMid}{\textbf{Бизнес}} & Сколько денег и довольных клиентов получает компания. Зависит от всего сразу, измерить трудно.\\
    \textcolor{acMid}{\textbf{Онлайн}} & Что видно на работающем сервисе: например, сколько времени люди сидят в приложении.\\
    \textcolor{acMid}{\textbf{Асессоры}} & Живые люди смотрят на ответы модели и ставят оценки. Нужно, когда качество --- дело вкуса.\\
    \textcolor{acMid}{\textbf{Офлайн}} & Считаем по старым данным до запуска: насколько ответы модели отличаются от правильных.\\
  \end{tabular}
\end{frame}

\begin{frame}[c]{Офлайн и онлайн метрики}
  \small
  \zebra{2}\small
  \begin{tabular}{>{\raggedright\arraybackslash}p{0.115\textwidth}>{\raggedright\arraybackslash}p{0.362\textwidth}>{\raggedright\arraybackslash}p{0.382\textwidth}}
    \rowcolor{acDark} & \textcolor{white}{\textbf{Офлайн}} & \textcolor{white}{\textbf{Онлайн}}\\
    Когда & До запуска модели & После деплоя, на работающей системе\\
    На чём & Исторические данные, тестовая выборка, оценки асессоров & Реальное поведение пользователей и магазинов\\
    Примеры & accuracy, MAE, MSE, precision, recall & Медианная длина сессии, число бананов на полках в конце дня\\
    Плюсы & Быстро, дёшево, без риска для пользователей & Ближе к бизнес-метрикам\\
    Минусы & Хороший результат не гарантирует рост бизнес-метрик & Нужны деплой, A/B-тест и риск для пользователей\\
  \end{tabular}
\end{frame}

\begin{frame}[c]{Путь модели от офлайна до бизнеса}
  \lead{Как всё связано.}

  \medskip
  \begin{columns}[T]
    \begin{column}{0.23\textwidth}
      \begin{block}{Офлайн}
        Метрики на тесте. Быстро отсеиваем слабые модели.
      \end{block}
    \end{column}
    \begin{column}{0.23\textwidth}
      \begin{block}{Деплой}
        Выкатка в продакшен. Нужна возможность отката.
      \end{block}
    \end{column}
    \begin{column}{0.23\textwidth}
      \begin{block}{Онлайн и A/B}
        Сравниваем с прошлой версией на случайных группах.
      \end{block}
    \end{column}
    \begin{column}{0.23\textwidth}
      \begin{block}{Бизнес}
        Доход и лояльность --- итоговая цель.
      \end{block}
    \end{column}
  \end{columns}

  \medskip
  \textbf{Офлайн-метрика --- прокси онлайн-метрики, онлайн --- прокси бизнеса.
  Хорошее число на тесте не гарантирует эффект в продакшене.}
\end{frame}

\begin{frame}[c]{Офлайн или онлайн?}
  \lead{Голосуем.}
  \begin{numrows}
        1 & Доля верно поставленных диагнозов на отложенной выборке\\
        2 & Среднее число бананов на полках в конце дня\\
        3 & Средний модуль разницы между предсказанной и истинной ценой квартиры\\
        4 & Медианная длина сессии в онлайн-игре\\
        5 & Оценка качества перевода асессорами\\
      \end{numrows}\par
\end{frame}

\begin{frame}[c]{Ответы: вид метрики}
  \small
  \zebra{2}
  \begin{tabular}{>{\raggedright\arraybackslash}p{0.324\textwidth}>{\raggedright\arraybackslash}p{0.553\textwidth}}
    \rowcolor{acDark} \textcolor{white}{\textbf{Метрика}} & \textcolor{white}{\textbf{Вид}}\\
    1. Доля верных диагнозов & Офлайн: считаем по историческим данным\\
    2. Бананы на полках & Онлайн: нужна работающая система\\
    3. Ошибка в цене квартиры & Офлайн: сравниваем с известными ценами\\
    4. Длина сессии & Онлайн: считается по реальным пользователям\\
    5. Оценка асессоров & Офлайн: оценивают до выпуска в продакшен\\
  \end{tabular}
\end{frame}

\begin{frame}[c]{Какую метрику выбрать?}
  \lead{Вопрос на подумать.}
  \small Модель отвечает «да» или «нет». Ошибиться можно двумя способами: сказать «да» зря
  (ложная тревога) или сказать «нет» зря (пропуск).

  \smallskip
  \begin{columns}[T]
    \begin{column}{0.31\textwidth}
      \begin{block}{Обычная поликлиника}
        Нужно, чтобы диагнозы в целом чаще были верными.
      \end{block}
    \end{column}
    \begin{column}{0.31\textwidth}
      \begin{block}{Проходная завода}
        Камера решает, сотрудник ли перед ней. Пустить чужого --- беда, а лишний раз не пустить своего --- мелочь.
      \end{block}
    \end{column}
    \begin{column}{0.31\textwidth}
      \begin{block}{Опасная болезнь}
        Пропущенный больной заразит других. Ложная тревога куда дешевле.
      \end{block}
    \end{column}
  \end{columns}

  \smallskip
  \textbf{Какую метрику выбрать в каждом случае?}
\end{frame}

\begin{frame}[c]{Ответ: разные нужды --- разные метрики}
  \small
  \zebra{1}
  \begin{tabular}{>{\raggedright\arraybackslash}p{0.191\textwidth}>{\raggedright\arraybackslash}p{0.688\textwidth}}
    \textcolor{acMid}{\textbf{Accuracy}}\newline{\scriptsize\color{acGray}обычная поликлиника} &
    \textbf{верные ответы $\div$ все ответы}\newline
    Какая доля ответов оказалась правильной. Подходит, когда обе ошибки одинаково неприятны.\\
    \textcolor{acMid}{\textbf{Precision}}\newline{\scriptsize\color{acGray}проходная завода} &
    \textbf{верно найденные $\div$ всё, что модель назвала «да»}\newline
    Можно ли верить модели, когда она говорит «да». Важна, когда ложное «да» дорого стоит: нельзя пустить чужого.\\
    \textcolor{acMid}{\textbf{Recall}}\newline{\scriptsize\color{acGray}опасная болезнь} &
    \textbf{верно найденные $\div$ все, кого нужно было найти}\newline
    Сколько из настоящих «да» модель нашла. Важна, когда нельзя пропустить: больной заразит других.\\
  \end{tabular}
\end{frame}

\begin{frame}[c]{Другие метрики}
  \lead{Вопрос на подумать.} Какую метрику выбрать здесь?

  \smallskip
  \begin{columns}[T]
    \begin{column}{0.31\textwidth}
      \begin{block}{Видеозвонок}
        Приложение «надевает» на участника виртуальный костюм рыцаря. Костюм должен ровно лечь на фигуру.
      \end{block}
    \end{column}
    \begin{column}{0.31\textwidth}
      \begin{block}{Безопасность}
        Охрана важного объекта: нужно заметить любого, кто вошёл в кадр.
      \end{block}
    \end{column}
    \begin{column}{0.31\textwidth}
      \begin{block}{Пандемия}
        Контроль, не превышает ли число людей в помещении порог.
      \end{block}
    \end{column}
  \end{columns}
\end{frame}

\begin{frame}[c]{Ответ: зависит от цели}
  \begin{columns}[T]
    \begin{column}{0.31\textwidth}
      \begin{block}{Точность положения}
        Среднеквадратичное отклонение координат опорных точек тела от истинных: костюм должен сесть ровно.
      \end{block}
    \end{column}
    \begin{column}{0.31\textwidth}
      \begin{block}{Полнота (recall)}
        На какой доле кадров с людьми модель их заметила. Точность рамки не так важна.
      \end{block}
    \end{column}
    \begin{column}{0.31\textwidth}
      \begin{block}{Число людей}
        Среднеквадратичное отклонение числа найденных людей от истинного.
      \end{block}
    \end{column}
  \end{columns}
\end{frame}

\begin{frame}[c]{Не всё сводится к метрикам}
  \begin{columns}[T]
    \begin{column}{0.31\textwidth}
      \begin{block}{Скорость}
        Предсказания в режиме реального времени. Требование и к модели, и к серверам.
      \end{block}
    \end{column}
    \begin{column}{0.31\textwidth}
      \begin{block}{Компактность}
        Модель должна помещаться на телефон или другое устройство.
      \end{block}
    \end{column}
    \begin{column}{0.31\textwidth}
      \begin{block}{Интерпретируемость}
        Можно объяснить решение: дадут ли кредит, согласуют ли лечение.
      \end{block}
    \end{column}
  \end{columns}
\end{frame}

\section{Данные}

\begin{frame}[c]{Данные: сколько и откуда}
  \small Чем сложнее задача, тем больше данных нужно. Они должны быть и многочисленными, и качественными.

  \smallskip
  \begin{columns}[T]
    \begin{column}{0.47\textwidth}
      \begin{block}{Открытые датасеты}
        Их много, но чаще для простых задач, например классификации картинок.
      \end{block}
      \begin{block}{Краудсорсинг}
        Люди размечают данные. Волонтёры без оплаты --- citizen science.
      \end{block}
    \end{column}
    \begin{column}{0.47\textwidth}
      \begin{block}{Покупка}
        Иногда датасет можно купить. Для некоторых задач данных нет нигде.
      \end{block}
      \begin{block}{Неразмеченные данные}
        Self-supervised learning и обучение представлений: используем данные без ответов.
      \end{block}
    \end{column}
  \end{columns}
\end{frame}

\begin{frame}[c]{Типы признаков}
  \small Признаки --- свойства объекта, по которым модель его описывает. Удобнее всего таблица: строки --- объекты, столбцы --- признаки.

  \smallskip
  \begin{columns}[T]
    \begin{column}{0.47\textwidth}
      \begin{block}{Численный}
        Рост, доход. Иногда отдельно выделяют вещественные и целочисленные.
      \end{block}
      \begin{block}{Бинарный}
        Два значения: да или нет, 0 или 1. Можно считать и числом, и категорией.
      \end{block}
    \end{column}
    \begin{column}{0.47\textwidth}
      \begin{block}{Категориальный}
        Значение из набора вариантов: профессия, день недели.
      \end{block}
      \begin{block}{Ординальный}
        Категории с порядком: класс опасности вещества от 1 до 4, год обучения.
      \end{block}
    \end{column}
  \end{columns}
\end{frame}

\begin{frame}[c]{Какого типа признак?}
  \lead{Вопрос на подумать.}
  \small Численный $\cdot$ Категориальный $\cdot$ Бинарный $\cdot$ Ординальный

  \smallskip
  \begin{columns}[T]
    \begin{column}{0.47\textwidth}
      \begin{numrows}
        1 & Рост человека\\
        2 & Профессия\\
        3 & Оплатил ли клиент заказ\\
      \end{numrows}\par
    \end{column}
    \begin{column}{0.47\textwidth}
      \begin{numrows}
        4 & Класс энергопотребления (от A до G)\\
        5 & День недели\\
        6 & Год обучения студента\\
      \end{numrows}\par
    \end{column}
  \end{columns}
\end{frame}

\begin{frame}[c]{Ответ: типы признаков}
  \small
  \zebra{2}
  \zebra{1}
  \begin{tabular}{>{\raggedright\arraybackslash}p{0.444\textwidth}l}
    \rowcolor{acDark} \textcolor{white}{\textbf{Признак}} & \textcolor{white}{\textbf{Тип}}\\
    Рост человека & \textcolor{acMid}{\textbf{Численный}}\\
    Профессия & \textcolor{acMid}{\textbf{Категориальный}}\\
    Оплатил ли заказ & \textcolor{acMid}{\textbf{Бинарный}}\\
    Класс энергопотребления, A--G & \textcolor{acMid}{\textbf{Ординальный}}\\
    День недели & \textcolor{acMid}{\textbf{Категориальный}}\\
    Год обучения & \textcolor{acMid}{\textbf{Ординальный}}\\
  \end{tabular}
\end{frame}

\begin{frame}[c]{Что не так с этими данными?}
  \lead{Найдите проблемы.}

  \medskip
  \large
  \zebra{2}\renewcommand{\arraystretch}{1.7}\setlength{\tabcolsep}{12pt}
  \begin{tabular}{llrl}
    \rowcolor{acDark} \textcolor{white}{\textbf{Клиент}} & \textcolor{white}{\textbf{Возраст}} & \textcolor{white}{\textbf{Доход, руб.}} & \textcolor{white}{\textbf{Вернул кредит}}\\
    Анна    & 34  & 90\,000    & да\\
    Борис   & N/A & 70\,000    & нет\\
    Виктор  & 140 & 85\,000    & да\\
    Галина  & 29  & 2\,400\,000 & да\\
    Дмитрий & 41  & 65\,000    & да\\
  \end{tabular}

  \bigskip
  \textbf{Часть проблем в таблице не видна вовсе. Какие это?}
\end{frame}

\begin{frame}[c]{Ответ: четыре типа проблем в данных}
  \begin{columns}[T]
    \begin{column}{0.47\textwidth}
      \begin{block}{Пропуски}
        Возраст Бориса неизвестен. Удалить, заполнить или оставить как признак: пропуск сам может нести информацию.
      \end{block}
      \begin{block}{Ошибки разметки}
        Разметчики ошибаются: бывают недобросовестные, а задача может быть сложной.
      \end{block}
    \end{column}
    \begin{column}{0.47\textwidth}
      \begin{block}{Выбросы}
        140 лет и доход 2,4 млн --- ошибка сбора или реальная аномалия. Важные выбросы обрабатывают отдельно.
      \end{block}
      \begin{block}{Data drift}
        Данные меняются со временем: формат или распределение. Нужен мониторинг.
      \end{block}
    \end{column}
  \end{columns}
\end{frame}

\section{Модель и обучение}

\begin{frame}[c]{Модель и алгоритм обучения}
  \begin{columns}[T]
    \begin{column}{0.47\textwidth}
      \begin{block}{Модель --- способ описать мир}
        «Земля плоская» --- тоже модель. В масштабе города она работает отлично,
        а на маршруте Париж --- Лос-Анджелес ломается, и её приходится заменить.
      \end{block}
    \end{column}
    \begin{column}{0.47\textwidth}
      \begin{block}{Предсказательная модель}
        \small
        \centerline{$f(x;\theta) \to y$}
        \smallskip
        $x$ --- признаки объекта, $y$ --- таргет, $\theta$ --- параметры, которые подбираются по данным.
      \end{block}
    \end{column}
  \end{columns}
\end{frame}

\begin{frame}[c]{Простейшая модель: одна константа}
  \small Всем квартирам предсказываем одну и ту же цену $c$, то есть $f(x)=c$.
  Цены пяти квартир, млн руб.: \textbf{5, 6, 7, 8, 24}.

  \smallskip
  \begin{columns}[T]
    \begin{column}{0.47\textwidth}
      \begin{block}{MAE}
        \centerline{$\displaystyle \mathrm{MAE}=\frac1N\sum_{i=1}^{N}\lvert f(x_i)-y_i\rvert$}
        \smallskip
        Средний модуль ошибки: насколько в среднем промахиваемся.

        \smallskip
        \lead{Лучшая $c$ --- медиана: 7}
      \end{block}
    \end{column}
    \begin{column}{0.47\textwidth}
      \begin{block}{MSE}
        \centerline{$\displaystyle \mathrm{MSE}=\frac1N\sum_{i=1}^{N}\bigl(f(x_i)-y_i\bigr)^2$}
        \smallskip
        Средний квадрат ошибки: большие промахи штрафуются сильнее.

        \smallskip
        \lead{Лучшая $c$ --- среднее: 10}
      \end{block}
    \end{column}
  \end{columns}

  \medskip
  {\footnotesize\color{acGray}
  Медиана --- число посередине, если выстроить цены по порядку. Одна дорогая квартира в 24 млн
  сдвинула среднее вверх, а медиану почти не тронула. Здесь $N$ --- число квартир, $i$ --- номер квартиры.\par}
\end{frame}

\begin{frame}[c]{Что нужно, чтобы обучить модель}
  \begin{columns}[T]
    \begin{column}{0.31\textwidth}
      \begin{block}{Модель}
        \small
        Формула с настраиваемыми числами $\theta$. Меняем $\theta$ --- меняются предсказания.
      \end{block}
    \end{column}
    \begin{column}{0.31\textwidth}
      \begin{block}{Функция потерь}
        \small
        Число, которое показывает, насколько модель ошибается. Чем меньше, тем лучше.
      \end{block}
    \end{column}
    \begin{column}{0.31\textwidth}
      \begin{block}{Алгоритм обучения}
        \small
        Подбирает $\theta$ так, чтобы потери стали меньше. Например, градиентный спуск:
        понемногу меняем $\theta$ туда, где ошибка убывает.
      \end{block}
    \end{column}
  \end{columns}

  \smallskip
  \begin{block}{Градиентный спуск в одну строку}
    \small
    \centerline{$\theta \leftarrow \theta-\eta\,\nabla L(\theta)$}
    \smallskip
    \small Здесь $\eta$ --- размер шага, $\nabla L$ --- направление, в котором ошибка растёт быстрее
    всего. Идём в противоположную сторону.
  \end{block}
\end{frame}

\begin{frame}[c]{Параметры и гиперпараметры}
  \begin{columns}[T]
    \begin{column}{0.47\textwidth}
      \begin{block}{Параметры}
        Числа внутри модели. Их подбирает алгоритм обучения по данным.
        \begin{itemize}
          \item коэффициенты $k$ и $b$ в прямой $y=kx+b$
          \item веса нейросети
          \item параметры $\theta$ из прошлого слайда
        \end{itemize}
      \end{block}
    \end{column}
    \begin{column}{0.47\textwidth}
      \begin{block}{Гиперпараметры}
        Настройки самого обучения. Их задаёт человек до начала обучения.
        \begin{itemize}
          \item шаг $\eta$ в градиентном спуске
          \item степень многочлена
          \item число слоёв нейросети
          \item сколько раз пройти по данным
        \end{itemize}
      \end{block}
    \end{column}
  \end{columns}
\end{frame}

\begin{frame}[c]{Недообучение и переобучение}
  \footnotesize\color{acGray}
  \textcolor{acOrange}{$\bullet$}~обучающие точки \qquad $\circ$~тестовые точки (новые данные)

  \smallskip
  \begin{columns}[T]
    \begin{column}{0.31\textwidth}
      \begin{block}{Слишком простая}
        \centering% underfit: straight line (degree 1).
% Least-squares fit (numpy.polyfit) to the 8 orange training points;
% hollow points are held-out test points. Data: y = 4(x-0.5)^2 + 0.15 + 0.25x
% + N(0, 0.2) noise, numpy seed 7.
% train MSE = 0.0992, test MSE (3 pts) = 0.0906
\begin{tikzpicture}[x=3.38cm,y=1.35cm]
  \draw[chart axis] (0,-0.12) -- (1.06,-0.12);
  \draw[chart axis] (0,-0.12) -- (0,1.42);
  \node[chart label,anchor=north east] at (1.06,-0.12) {$x$};
  \node[chart label,anchor=north east] at (0,1.42) {$y$};
  \begin{scope}
    \clip (0.01,-0.10) rectangle (1.04,1.40);
    \draw[chart curve] plot[smooth] coordinates {(0.025,0.401) (0.035,0.404) (0.044,0.408) (0.054,0.412) (0.063,0.416) (0.073,0.420) (0.083,0.424) (0.092,0.428) (0.102,0.432) (0.111,0.436) (0.121,0.440) (0.131,0.443) (0.140,0.447) (0.150,0.451) (0.159,0.455) (0.169,0.459) (0.179,0.463) (0.188,0.467) (0.198,0.471) (0.207,0.475) (0.217,0.478) (0.227,0.482) (0.236,0.486) (0.246,0.490) (0.255,0.494) (0.265,0.498) (0.275,0.502) (0.284,0.506) (0.294,0.510) (0.303,0.513) (0.313,0.517) (0.323,0.521) (0.332,0.525) (0.342,0.529) (0.351,0.533) (0.361,0.537) (0.371,0.541) (0.380,0.545) (0.390,0.549) (0.399,0.552) (0.409,0.556) (0.419,0.560) (0.428,0.564) (0.438,0.568) (0.447,0.572) (0.457,0.576) (0.467,0.580) (0.476,0.584) (0.486,0.587) (0.495,0.591) (0.505,0.595) (0.515,0.599) (0.524,0.603) (0.534,0.607) (0.543,0.611) (0.553,0.615) (0.563,0.619) (0.572,0.623) (0.582,0.626) (0.591,0.630) (0.601,0.634) (0.611,0.638) (0.620,0.642) (0.630,0.646) (0.639,0.650) (0.649,0.654) (0.659,0.658) (0.668,0.661) (0.678,0.665) (0.687,0.669) (0.697,0.673) (0.707,0.677) (0.716,0.681) (0.726,0.685) (0.735,0.689) (0.745,0.693) (0.755,0.696) (0.764,0.700) (0.774,0.704) (0.783,0.708) (0.793,0.712) (0.803,0.716) (0.812,0.720) (0.822,0.724) (0.831,0.728) (0.841,0.732) (0.851,0.735) (0.860,0.739) (0.870,0.743) (0.879,0.747) (0.889,0.751) (0.899,0.755) (0.908,0.759) (0.918,0.763) (0.927,0.767) (0.937,0.770) (0.947,0.774) (0.956,0.778) (0.966,0.782) (0.975,0.786) (0.985,0.790)};
  \end{scope}
  \node[chart test] at (0.230,0.230) {};
  \node[chart test] at (0.480,0.180) {};
  \node[chart test] at (0.870,0.535) {};
  \node[chart train] at (0.065,0.825) {};
  \node[chart train] at (0.176,0.490) {};
  \node[chart train] at (0.311,0.469) {};
  \node[chart train] at (0.409,0.357) {};
  \node[chart train] at (0.542,0.314) {};
  \node[chart train] at (0.675,0.255) {};
  \node[chart train] at (0.780,0.653) {};
  \node[chart train] at (0.953,1.348) {};
\end{tikzpicture}%
\par
        Прямая не ловит изгиб данных.
      \end{block}
    \end{column}
    \begin{column}{0.31\textwidth}
      \begin{block}{Нормальная}
        \centering% good fit: parabola (degree 2).
% Least-squares fit (numpy.polyfit) to the 8 orange training points;
% hollow points are held-out test points. Data: y = 4(x-0.5)^2 + 0.15 + 0.25x
% + N(0, 0.2) noise, numpy seed 7.
% train MSE = 0.0102, test MSE (3 pts) = 0.0780
\begin{tikzpicture}[x=3.38cm,y=1.35cm]
  \draw[chart axis] (0,-0.12) -- (1.06,-0.12);
  \draw[chart axis] (0,-0.12) -- (0,1.42);
  \node[chart label,anchor=north east] at (1.06,-0.12) {$x$};
  \node[chart label,anchor=north east] at (0,1.42) {$y$};
  \begin{scope}
    \clip (0.01,-0.10) rectangle (1.04,1.40);
    \draw[chart curve] plot[smooth] coordinates {(0.025,0.991) (0.035,0.958) (0.044,0.926) (0.054,0.895) (0.063,0.864) (0.073,0.835) (0.083,0.806) (0.092,0.777) (0.102,0.750) (0.111,0.723) (0.121,0.697) (0.131,0.672) (0.140,0.647) (0.150,0.624) (0.159,0.601) (0.169,0.578) (0.179,0.557) (0.188,0.536) (0.198,0.516) (0.207,0.496) (0.217,0.478) (0.227,0.460) (0.236,0.443) (0.246,0.427) (0.255,0.411) (0.265,0.396) (0.275,0.382) (0.284,0.369) (0.294,0.356) (0.303,0.344) (0.313,0.333) (0.323,0.322) (0.332,0.313) (0.342,0.304) (0.351,0.296) (0.361,0.288) (0.371,0.281) (0.380,0.275) (0.390,0.270) (0.399,0.266) (0.409,0.262) (0.419,0.259) (0.428,0.257) (0.438,0.255) (0.447,0.254) (0.457,0.254) (0.467,0.255) (0.476,0.257) (0.486,0.259) (0.495,0.262) (0.505,0.265) (0.515,0.270) (0.524,0.275) (0.534,0.281) (0.543,0.288) (0.553,0.295) (0.563,0.303) (0.572,0.312) (0.582,0.322) (0.591,0.332) (0.601,0.343) (0.611,0.355) (0.620,0.368) (0.630,0.381) (0.639,0.395) (0.649,0.410) (0.659,0.425) (0.668,0.442) (0.678,0.459) (0.687,0.477) (0.697,0.495) (0.707,0.514) (0.716,0.535) (0.726,0.555) (0.735,0.577) (0.745,0.599) (0.755,0.622) (0.764,0.646) (0.774,0.670) (0.783,0.696) (0.793,0.721) (0.803,0.748) (0.812,0.776) (0.822,0.804) (0.831,0.833) (0.841,0.862) (0.851,0.893) (0.860,0.924) (0.870,0.956) (0.879,0.989) (0.889,1.022) (0.899,1.056) (0.908,1.091) (0.918,1.127) (0.927,1.163) (0.937,1.200) (0.947,1.238) (0.956,1.277) (0.966,1.316) (0.975,1.356) (0.985,1.397)};
  \end{scope}
  \node[chart test] at (0.230,0.230) {};
  \node[chart test] at (0.480,0.180) {};
  \node[chart test] at (0.870,0.535) {};
  \node[chart train] at (0.065,0.825) {};
  \node[chart train] at (0.176,0.490) {};
  \node[chart train] at (0.311,0.469) {};
  \node[chart train] at (0.409,0.357) {};
  \node[chart train] at (0.542,0.314) {};
  \node[chart train] at (0.675,0.255) {};
  \node[chart train] at (0.780,0.653) {};
  \node[chart train] at (0.953,1.348) {};
\end{tikzpicture}%
\par
        Улавливает общий тренд, не гонясь за каждой точкой.
      \end{block}
    \end{column}
    \begin{column}{0.31\textwidth}
      \begin{block}{Переобученная}
        \centering% overfit: polynomial of degree 7.
% Least-squares fit (numpy.polyfit) to the 8 orange training points;
% hollow points are held-out test points. Data: y = 4(x-0.5)^2 + 0.15 + 0.25x
% + N(0, 0.2) noise, numpy seed 7.
% train MSE = 0.0000, test MSE (3 pts) = 0.3830
\begin{tikzpicture}[x=3.38cm,y=1.35cm]
  \draw[chart axis] (0,-0.12) -- (1.06,-0.12);
  \draw[chart axis] (0,-0.12) -- (0,1.42);
  \node[chart label,anchor=north east] at (1.06,-0.12) {$x$};
  \node[chart label,anchor=north east] at (0,1.42) {$y$};
  \begin{scope}
    \clip (0.01,-0.10) rectangle (1.04,1.40);
    \draw[chart curve] plot[smooth] coordinates {(0.025,1.870) (0.035,1.523) (0.044,1.244) (0.054,1.023) (0.063,0.849) (0.073,0.717) (0.083,0.619) (0.092,0.548) (0.102,0.500) (0.111,0.470) (0.121,0.454) (0.131,0.449) (0.140,0.451) (0.150,0.458) (0.159,0.469) (0.169,0.481) (0.179,0.494) (0.188,0.505) (0.198,0.516) (0.207,0.524) (0.217,0.530) (0.227,0.533) (0.236,0.533) (0.246,0.532) (0.255,0.527) (0.265,0.521) (0.275,0.512) (0.284,0.502) (0.294,0.491) (0.303,0.479) (0.313,0.466) (0.323,0.453) (0.332,0.439) (0.342,0.426) (0.351,0.414) (0.361,0.402) (0.371,0.391) (0.380,0.381) (0.390,0.372) (0.399,0.364) (0.409,0.357) (0.419,0.351) (0.428,0.346) (0.438,0.342) (0.447,0.338) (0.457,0.336) (0.467,0.333) (0.476,0.331) (0.486,0.330) (0.495,0.328) (0.505,0.326) (0.515,0.323) (0.524,0.320) (0.534,0.317) (0.543,0.313) (0.553,0.308) (0.563,0.303) (0.572,0.297) (0.582,0.291) (0.591,0.284) (0.601,0.277) (0.611,0.270) (0.620,0.264) (0.630,0.258) (0.639,0.254) (0.649,0.251) (0.659,0.250) (0.668,0.252) (0.678,0.257) (0.687,0.266) (0.697,0.279) (0.707,0.297) (0.716,0.321) (0.726,0.350) (0.735,0.387) (0.745,0.430) (0.755,0.481) (0.764,0.539) (0.774,0.605) (0.783,0.679) (0.793,0.760) (0.803,0.847) (0.812,0.942) (0.822,1.041) (0.831,1.144) (0.841,1.248) (0.851,1.353) (0.860,1.455) (0.870,1.550) (0.879,1.636) (0.889,1.707) (0.899,1.758) (0.908,1.784) (0.918,1.778) (0.927,1.732) (0.937,1.637) (0.947,1.484) (0.956,1.261) (0.966,0.957) (0.975,0.559) (0.985,0.050)};
  \end{scope}
  \node[chart test] at (0.230,0.230) {};
  \node[chart test] at (0.480,0.180) {};
  \node[chart test] at (0.870,0.535) {};
  \node[chart train] at (0.065,0.825) {};
  \node[chart train] at (0.176,0.490) {};
  \node[chart train] at (0.311,0.469) {};
  \node[chart train] at (0.409,0.357) {};
  \node[chart train] at (0.542,0.314) {};
  \node[chart train] at (0.675,0.255) {};
  \node[chart train] at (0.780,0.653) {};
  \node[chart train] at (0.953,1.348) {};
\end{tikzpicture}%
\par
        Проходит почти через каждую обучающую точку, но между ними и на новых данных ошибается.
      \end{block}
    \end{column}
  \end{columns}
\end{frame}

\begin{frame}[c]{Ошибка на train и на test}
  \begin{columns}[T]
    \begin{column}{0.56\textwidth}
      \centering
      % Train/test MSE vs polynomial degree 0..5 for the same data as fit-*.tex.
% Train MSE on the 8 training points, test MSE on 400 held-out points
% (numpy.polyfit; seed 7, noise 0.2).
% train: 0.1126, 0.0992, 0.0102, 0.0050, 0.0041, 0.0011
% test:  0.1082, 0.1029, 0.0356, 0.0380, 0.0393, 0.0477
\begin{tikzpicture}[x=1.06cm,y=1cm]
  \draw[chart grid] (0,0.000) -- (5.25,0.000);
  \node[chart label,anchor=east] at (0,0.000) {0{,}000};
  \draw[chart grid] (0,0.570) -- (5.25,0.570);
  \node[chart label,anchor=east] at (0,0.570) {0{,}025};
  \draw[chart grid] (0,1.140) -- (5.25,1.140);
  \node[chart label,anchor=east] at (0,1.140) {0{,}050};
  \draw[chart grid] (0,1.710) -- (5.25,1.710);
  \node[chart label,anchor=east] at (0,1.710) {0{,}075};
  \draw[chart grid] (0,2.280) -- (5.25,2.280);
  \node[chart label,anchor=east] at (0,2.280) {0{,}100};
  \draw[chart grid] (0,2.850) -- (5.25,2.850);
  \node[chart label,anchor=east] at (0,2.850) {0{,}125};
  \draw[chart axis] (0,0) -- (5.25,0);
  \draw[chart axis] (0,0) -- (0,3.000);
  \foreach \d in {0,...,5} {\draw[chart axis] (\d,0) -- (\d,-0.07); \node[chart label,anchor=north] at (\d,-0.07) {\d};}
  \draw[acAccent!60, line width=0.6pt, dash pattern=on 2pt off 2pt] (2,0) -- (2,0.812);
  \draw[acInk, line width=1.3pt, line join=round] plot coordinates {(0,2.567) (1,2.261) (2,0.233) (3,0.114) (4,0.093) (5,0.024)};
  \draw[acOrange, line width=1.3pt, line join=round] plot coordinates {(0,2.467) (1,2.346) (2,0.812) (3,0.867) (4,0.895) (5,1.087)};
  \fill[acInk] (0,2.567) circle (1.9pt);
  \fill[acInk] (1,2.261) circle (1.9pt);
  \fill[acInk] (2,0.233) circle (1.9pt);
  \fill[acInk] (3,0.114) circle (1.9pt);
  \fill[acInk] (4,0.093) circle (1.9pt);
  \fill[acInk] (5,0.024) circle (1.9pt);
  \fill[acOrange] (0,2.467) circle (1.9pt);
  \fill[acOrange] (1,2.346) circle (1.9pt);
  \fill[acOrange] (2,0.812) circle (1.9pt);
  \fill[acOrange] (3,0.867) circle (1.9pt);
  \fill[acOrange] (4,0.895) circle (1.9pt);
  \fill[acOrange] (5,1.087) circle (1.9pt);
  \draw[acAccent, line width=1.2pt] (2,0.812) circle (5.5pt);
  \node[font=\scriptsize\bfseries, text=acAccent, anchor=east, inner sep=1pt] at (2-0.17,0.812) {оптимум};
  % legend
  \begin{scope}[shift={(3.35,2.800)}]
    \fill[acCard] (-0.1,0.2) rectangle (1.15,-0.56);
    \draw[acInk, line width=1.3pt] (0,0) -- (0.4,0); \fill[acInk] (0.2,0) circle (1.9pt);
    \node[chart label, text=acText, anchor=west] at (0.48,0) {train};
    \draw[acOrange, line width=1.3pt] (0,-0.36) -- (0.4,-0.36); \fill[acOrange] (0.2,-0.36) circle (1.9pt);
    \node[chart label, text=acText, anchor=west] at (0.48,-0.36) {test};
  \end{scope}
  \node[chart label, anchor=north] at (2.625,-0.42) {Степень многочлена (сложность модели)};
  \node[chart label, anchor=south, rotate=90] at (-0.78,1.425) {MSE};
\end{tikzpicture}%

    \end{column}
    \begin{column}{0.40\textwidth}
      \small
      \textbf{train} (тёмная линия) падает с ростом сложности, иногда до нуля.

      \smallskip
      \textbf{test} (оранжевая линия) сначала падает, потом снова растёт.

      \smallskip
      \lead{Оптимум --- модель не слишком простая и не слишком сложная.}

      \smallskip
      \footnotesize\color{acGray}Метрики считаем на тесте и никогда не смешиваем его с обучающей выборкой.
      По горизонтали растёт сложность модели (степень многочлена от 0 до 5), по вертикали --- MSE.
    \end{column}
  \end{columns}
\end{frame}

\begin{frame}[c]{После обучения}
  \begin{columns}[T]
    \begin{column}{0.47\textwidth}
      \begin{block}{Деплой}
        Выкатка в продакшен: эффективный код, параллельность, дружба с фреймворками.
      \end{block}
      \begin{block}{Мониторинг}
        Следим за качеством, дообучаем и переобучаем на новых данных.
      \end{block}
    \end{column}
    \begin{column}{0.47\textwidth}
      \begin{block}{A/B-тест}
        Сравнение с предыдущей версией на случайных группах. Если хуже --- откат.
      \end{block}
      \begin{block}{Drift}
        Data drift: меняются данные. Concept drift: меняется связь признаков и таргета.
      \end{block}
    \end{column}
  \end{columns}
\end{frame}


\end{document}
